% =========================
% Cas pratique 2
% =========================

\section[Cas 2]{Cas pratique 2 : La saisie-contrefaçon}

%--------------------------
\begin{frame}{Question 1}

\begin{block}{}
Art. \textbf{D211-6} COJ dispose que le tribunal judiciaire ayant compétence exclusive pour connaître des actions en matière de brevets d'invention, de certificats d'utilité, de certificats complémentaires de protection et de topographies de produits semi-conducteurs, dans les cas et conditions prévus par le code de la propriété intellectuelle, est celui de Paris.
\end{block}

\begin{itemize}
  \item En l'espèce, la saisie-contrefaçon en brevet est ordonnée par le Tribunal judiciaire de Paris.
  De plus, France Tube se prévaut d'une action en contrefaçon de la partie française d'un brevet européen.
  \item En conclusion, le tribunal jurdiciaire de Paris est compétent pour ordonner la saisie-contrefaçon.
\end{itemize}

\end{frame}
%--------------------------

%--------------------------
\begin{frame}{Question 2}

\begin{block}{}
Art. \textbf{L615-5} al. 3 CPI dispose que la juridiction peut ordonner, aux mêmes fins probatoires, la description détaillée ou la saisie réelle 
[...]
\end{block}

\begin{itemize}
  \item En l'espèce, la saisie intégrale du stock n'est pas une saisie probatoire.
  \item En conclusion, la saisie du stock ne semble pas licite.
\end{itemize}

\end{frame}

\begin{frame}{Question 2}
    
\begin{block}{}
Art. \textbf{496} CPC dispose que s'il est fait droit à la requête, tout intéressé peut en référer au juge qui a rendu l'ordonnance.
\end{block}
\begin{itemize}
    \item En l'espèce, le juge a fait droit à la requête de saisie-contrefaçon.
     Mais, la saisie ordonnée présente un élément illicite.
    \item En conclusion, la constestation se fait par la procédure de référé.
\end{itemize}
\begin{block}{}
Art. \textbf{497} CPC dispose que le juge a la faculté de modifier ou de rétracter son ordonnance, même si le juge du fond est saisi de l'affaire.
\end{block}
\begin{itemize}
    \item En l'espèce, la contestation semble légitime.
    Le juge a un motif pour rétracter l'ordonnance.
    \item En conclusion, le pv de la saisie est susceptible d'être annulé. 
\end{itemize}
    
\end{frame}
%--------------------------

%--------------------------
\begin{frame}{Question 3}
    
\begin{block}{}
\textbf{Principe de loyauté} (CA 06/04/2022 7ème page 2e par.): "ne donnant pas au juge des requêtes
tous les éléments en sa possession pour qu'il puisse exercer son contrôle".
\end{block}

\begin{itemize}
    \item En l'espèce, l'information sur l'opposition est pertinente pour ordonner la saisie-contrefaçon.
    Il y a en effet un sursit à statuer dans ce cas-là (Art. 378 CPC).
    \item En conclusion, l'ordonnance est susceptible d'être rétractée.
\end{itemize}

\end{frame}
%--------------------------

%--------------------------
\begin{frame}{Question 4}

\begin{block}{}
L'art. \textbf{60} AJUB dispose que la Juridication peut, avant même
l'engagement d'une action au fond, ordonner des mesures provisoires rapides
et efficaces pour conserver les éléments de preuve pertinents au regard
de la contrefaçon alléguée [...]
\end{block}

\begin{itemize}
  \item En l'espèce, France Tube a la possibilité de solliciter devant la JUB une mesure de préservation de preuve.
  \item En conclusion, France Tube a à sa disposition une mesure équivalente à la saisie-contrefaçon française.
\end{itemize}

\end{frame}

%--------------------------

\begin{frame}{Question 4}
  
\begin{block}{}
L'art. \textbf{60} AJUB dispose que à la demande du requérant
qui a présenté des éléments de preuve raisonnablement accessibles
pour étayer ses allégations selon lesquelles sont brevet a
été contrefait ou qu'une contrefaçon est imminente [...].
\end{block}

\begin{itemize}
  \item En l'espèce, il y a trois procès-verbal.
  \item En conclusion, une condition pour requérir à une mesure de préservation de preuve
  est remplie.
\end{itemize}

  \begin{table}[h]
\centering
\begin{tabular}{l l}
  Saisie CF JUB & Saisie CF FR \\
  \hline
  + portée territoriale & - portée territoriale\\
  procédure contradictoire & non contradictoire \\
  motifs requis (R192 RPJUB) & pas de preuve requise \\
\end{tabular}
\end{table}
\end{frame}